\chapter{Impact\\แปลผลงานเป็นการเปลี่ยนแปลงที่มีความหมาย}

Output คือสิ่งที่โครงการผลิต เช่น บทความ ต้นแบบ หรือการอบรม ส่วน Impact คือความเปลี่ยนแปลงที่เกิดแก่ผู้คน ระบบ เศรษฐกิจ สิ่งแวดล้อม หรือนโยบาย การแยกสองสิ่งนี้ช่วยให้ Mentee ไม่หยุดที่การส่งมอบผลงาน แต่ติดตามว่าผลงานเดินทางไปสร้างคุณค่าอย่างไร

\learningobjectives{
  \item แยก output, outcome และ impact
  \item เขียน Impact Statement ที่ระบุผู้รับคุณค่าและหลักฐาน
  \item หลีกเลี่ยงการกล่าวอ้างเกินส่วนของโครงการ
}

\section{เส้นทางจากงานสู่คุณค่า}
\begin{longtable}{P{.18\textwidth}P{.34\textwidth}P{.34\textwidth}}
\toprule
\textbf{ระดับ} & \textbf{คำถาม} & \textbf{ตัวอย่าง}\\
\midrule
Output & เราส่งมอบอะไร & ระบบคัดแยกและคู่มือใช้งาน\\
Outcome & ใครเปลี่ยนพฤติกรรมหรือผลการทำงานอย่างไร & ร้านค้าใช้ระบบและลดเวลาคัดแยก\\
Impact & การเปลี่ยนแปลงระยะยาวหรือระดับระบบคืออะไร & ของเสียลดลง รายได้สุทธิเพิ่ม และเกิดแนวปฏิบัติใหม่\\
\bottomrule
\end{longtable}

\section{Impact Statement}
Impact Statement ที่ดีระบุ \textbf{ผู้ที่เปลี่ยนแปลง} \textbf{สิ่งที่เปลี่ยน} \textbf{ขนาดหรือคุณภาพของการเปลี่ยน} \textbf{ช่วงเวลา} และ \textbf{บทบาทของงาน} เช่น \enquote{ภายในหกเดือน ร้านค้านำร่อง 12 แห่งลดอาหารเสียเฉลี่ยร้อยละ 45 และประหยัดต้นทุนรวมประมาณ 500,000 บาทต่อปี โดยใช้ระบบคัดแยกและช่องทางขายที่ทีมพัฒนาร่วมกับผู้ประกอบการ}

\section{Contribution ไม่ใช่ Attribution ทั้งหมด}
Impact ส่วนใหญ่เกิดจากหลายปัจจัย Mentee ควรอธิบายอย่างระมัดระวังว่างาน \enquote{มีส่วนสนับสนุน} อย่างไร แทนการอ้างว่าเป็นสาเหตุทั้งหมด Logic model หรือ theory of change ช่วยเปิดสมมติฐานระหว่างกิจกรรมกับผลกระทบ และชี้ว่าต้องมีเงื่อนไขใดร่วมด้วย

% Figure placeholder: Output-to-impact pathway

\begin{examplebox}
บทความได้รับการอ้างอิง 80 ครั้งเป็นสัญญาณการใช้ความรู้ แต่ยังไม่ใช่ Impact สาธารณะโดยตัวมันเอง เมื่อพบว่าหน่วยงานนำวิธีจากบทความไปปรับมาตรฐานตรวจคุณภาพสำหรับสหกรณ์ 24 แห่ง จึงเริ่มเห็นเส้นทางจากความรู้สู่การเปลี่ยนแปลงระดับระบบ
\end{examplebox}

\diagnostic{หากผลงานนี้ไม่เกิดขึ้น ใครจะเผชิญสถานการณ์ต่างไปอย่างไร และมีข้อมูลใดยืนยัน}

\begin{mentornote}
Impact ที่เล็กแต่ยืนยันได้มีคุณค่ากว่า Impact ใหญ่ที่เป็นเพียงความหวัง ช่วย Mentee แยก \enquote{ผลที่เกิดแล้ว} \enquote{ผลที่กำลังติดตาม} และ \enquote{ผลที่คาดหวัง}
\end{mentornote}

\begin{keytakeaway}
Impact ทำให้คุณค่าของงานอ่านออกสำหรับผู้มีส่วนได้ส่วนเสีย โดยเชื่อม output กับความเปลี่ยนแปลงจริงอย่างมีหลักฐานและไม่กล่าวอ้างเกินส่วน
\end{keytakeaway}

\begin{reflection}
เลือก output หนึ่งชิ้นและเขียนเส้นทางไปสู่ outcome และ impact ระบุสมมติฐานสำคัญสองข้อ ผู้มีส่วนได้ส่วนเสียสามกลุ่ม และตัวชี้วัดหนึ่งตัวต่อกลุ่ม
\blankline\blankline
\end{reflection}

\chaptersummary{การสื่อสาร Impact ต้องเริ่มจากผู้รับคุณค่าและความเปลี่ยนแปลง มิใช่เริ่มจากผลงานเพียงอย่างเดียว Mentor ช่วย Mentee สร้างเส้นทางเชิงเหตุผล เก็บหลักฐาน และสื่อสาร contribution อย่างซื่อตรง}
